%------------------------------------------------------------------------------
% Início do abstract
%------------------------------------------------------------------------------
\begin{center}
    \textbf{ABSTRACT}
    \vspace{10pt}
\end{center}

This work investigates the use of large language models (LLMs) as an instrument to support scientific writing in the context of Brazilian conferences. The growth of scientific production has intensified structural limitations of the peer review process, such as reviewer overload, heterogeneity of reviews, and scarcity of structured feedback for early-career authors. A particularly burdensome consequence for new researchers is early editorial rejection (\textit{desk rejection}): the refusal of a manuscript before merit review, frequently motivated not by scientific flaws, but by the inadequacy of the text to the tacit conventions of form and content expected by a specific track or conference. For this purpose, the design of a software artifact capable of acting as a scientific writing tutor is investigated. The proposal repositions automated review: instead of evaluating the manuscript against universal quality criteria, the artifact first models the communicative genre of a target track from the corpus of already accepted papers and then verifies the compliance of the submitted manuscript to this specification, generating actionable diagnostics.
The theoretical foundation articulates Genre Theory applied to Information Systems and the analysis of rhetorical \textit{moves}, in the formulation of the CARS model, integrated under the Design Science Research (DSR) method, involving literature review, artifact development and empirical evaluation through a case study with researchers.
An efficient method for analyzing the set of papers of a conference is also described: structured extraction per paper in a local model, followed by deterministic aggregation, minimizing computational cost.
It is expected to contribute with a usable artifact, a set of generalizable design principles and empirical evidence on the reduction of early rejection and support for the training of researchers in the Brazilian context.
\vspace{20pt}

\textbf{Keywords:}~Scientific Writing; Genre Analysis; Rhetorical \textit{Moves}; Large Language Models (LLM); Design Science Research;

%------------------------------------------------------------------------------
% Fim do abstract
%------------------------------------------------------------------------------
